\documentclass[../../../include/open-logic-section]{subfiles}

\begin{document}

\olfileid{sth}{story}{approach}
\olsection{رویکرد انباشتی-تکراری}

در ادامه شرحی اندکی کامل‌تر از چگونگی لایه‌بندی مجموعه‌ها در مراحل می‌آید:
\begin{quote}
	مجموعه‌ها در \emph{مراحل} تشکیل می‌شوند. برای هر مرحلهٔ~$S$، مراحل
	معینی وجود دارند که \emph{پیش از}~$S$ هستند. در مرحلهٔ~$S$، هر گردایه
	متشکل از مجموعه‌هایی که در مراحل پیش از~$S$ تشکیل شده‌اند، به صورت یک
	مجموعه تشکیل می‌شود. هیچ مجموعه‌ای جز مجموعه‌هایی که در مراحل تشکیل
	می‌شوند وجود ندارد. \citep[p.~323]{Shoenfield:AST}
\end{quote}
این طرحی از \emph{تصور انباشتی-تکراریِ مجموعه} است. این تصور زیربنای
نظریهٔ صوری مجموعه‌ها خواهد بود که در~\olref[sth][][]{part} عرضه می‌کنیم.

بیایید اندکی دقیق‌تر آن را بررسی کنیم. در فرایندی که شوئنفیلد وصف می‌کند،
در هر مرحله مجموعه‌های تازه‌ای را از مجموعه‌هایی می‌سازیم که از مراحل
پیشین در دسترس ما بوده‌اند. بنابراین، در تصویر شوئنفیلد، در مرحلهٔ آغازین،
یعنی مرحلهٔ~$0$، هیچ مرحلهٔ \emph{پیشینی} وجود ندارد و در نتیجه، به طریق
اولی، هیچ مجموعه‌ای از مراحل پیشین در دسترس ما نیست.\footnote{چرا باید
فرض کنیم که مرحلهٔ نخستینی \emph{وجود دارد}؟ بنگرید به پانوشتِ
\stagesord{} در~\olref[z][story]{sec}.} پس تنها یک مجموعه می‌سازیم:
مجموعهٔ بدون هیچ !!{element}، یعنی~$\emptyset$. در مرحلهٔ~$1$، دقیقاً یک
مجموعه از مراحل پیشین در دسترس ماست، پس تنها مجموعهٔ تازه
$\{\emptyset\}$ است. در مرحلهٔ~$2$، دو مجموعه از مراحل پیشین در دسترس
ماست و دو مجموعهٔ تازهٔ~$\{\{\emptyset\}\}$ و
$\{\emptyset, \{\emptyset\}\}$ را می‌سازیم. در مرحلهٔ~$3$، چهار مجموعه
از مراحل پیشین در دسترس ماست و دوازده مجموعهٔ تازه می‌سازیم\ldots. بدین
ترتیب، تصویر انباشتی-تکراریِ مجموعه‌ها کم‌وبیش چنین خواهد بود (عددها مراحل
را نشان می‌دهند):
\begin{center}
	\begin{tikzpicture}[scale=0.6]
	\tikzset{cut_here/.style={densely dotted}}
	\draw (-4,6) -- (-1, 0)-- (2,6);
	\draw[cut_here] (-5,8)--(-4,6);
	\draw[cut_here] (2,6)--(3,8);
	\draw(-1.5, 1)--(-0.5, 1);
	\draw(-2, 2)--(0, 2);
	\draw(-2.5, 3)--(0.5,3);
	\draw(-3, 4)--(1, 4);
	\draw(-3.5, 5)--(1.5, 5);
	\draw(-4, 6)--(2, 6);
	\node[label] at (0, 0) {\small 0};
	\node[label] at (0.5, 1) {\small 1};
	\node[label] at (1, 2) {\small 2};
	\node[label] at (1.5, 3) {\small 3};
	\node[label] at (2, 4) {\small 4};
	\node[label] at (2.5, 5) {\small 5};
	\node[label] at (3, 6) {\small 6};
	\end{tikzpicture}
\end{center}
پس چرا باید این روایت را بپذیریم؟

یک دلیل آن است که روایتی خوش‌ساخت و مهارپذیر است. اکنون که آشکارترین
روایت، یعنی تصریح ساده‌انگارانه، شکست خورده است، به روایتی مناسب نیاز داریم.

اما این روایت \emph{فقط} خوش‌ساخت نیست. دلیل خوبی داریم که باور کنیم هر
نظریهٔ مجموعه‌هایی که بر این روایت بنا شود \emph{سازگار} خواهد بود. دلیل
چنین است.

بنا بر تصور انباشتی-تکراریِ مجموعه، مجموعه‌ها را در مراحل می‌سازیم و
!!{element}های آن‌ها باید اشیایی باشند که \emph{از پیش} در دسترس بوده‌اند.
پس برای هر مرحلهٔ~$S$ می‌توانیم مجموعهٔ زیر را بسازیم:
\[
	R_S = \Setabs{x}{x \notin x \text{ و $x$ پیش از $S$ در دسترس بوده است}}
\]
استدلالی که پارادوکس راسل را اثبات می‌کند، اکنون نشان خواهد داد که خود
$R_S$ پیش از مرحلهٔ~$S$ در دسترس نیست. و این تناقضی دربر ندارد. افزون بر
این، اگر تصور انباشتی-تکراریِ مجموعه را بپذیریم، حتی نباید \emph{انتظار}
داشته باشیم بتوانیم خود مجموعهٔ راسل را بسازیم؛ زیرا آن مجموعهٔ همهٔ
مجموعه‌های خودناعضوی خواهد بود که «در هر زمانی در دسترس قرار خواهند
گرفت». خلاصه اینکه، این واقعیت که (به‌طور اثبات‌پذیر) نمی‌توانیم مجموعهٔ
راسل را بسازیم، بنا بر روایت انباشتی-تکراری \emph{شگفت‌آور} نیست؛ درست
همان چیزی است که \emph{پیش‌بینی} می‌کنیم.

%پس تصور انباشتی-تکراریِ مجموعه از یک جهت پاسخی به پارادوکس راسل می‌دهد که بسیار شبیه پاسخ \emph{محمولی‌گرا} است. زیرا محمولی‌گرا می‌گفت گردایهٔ همهٔ مجموعه‌های$_n$ خودناعضو، مجموعه‌ای$_{n+1}$ است و این مجموعه$_n$ خودعضو نخواهد بود. (در واقع، بیشتر محمولی‌گرایان پرسش از خودعضوبودن یک مجموعه را \emph{به‌لحاظ دستوری نادرست} می‌شمارند.)
%
%اما میان رویکرد انباشتی-تکراری و رویکرد محمولی‌گرا تفاوت مهمی وجود دارد: رویکرد انباشتی-تکراری همهٔ موجودات ما را \emph{از یک نوع} می‌داند. محمولی‌گرا احتمالاً یک مجموعهٔ تهی$_0$ (یعنی مجموعه‌ای$_0$ بدون هیچ !!{element}) و یک مجموعهٔ تهی$_1$ (یعنی مجموعه‌ای$_1$ بدون هیچ !!{element}) خواهد داشت و این دو \emph{موجوداتی متفاوت} خواهند بود. اما در رویکرد انباشتی-تکراری تنها یک مجموعهٔ تهی، یعنی~$\emptyset$، وجود دارد و آن در هر مرحلهٔ سلسله‌مراتب در دسترس خواهد بود.

%در واقع، اصل موضوع گسترش‌مندی که در آغاز این فصل بیان شد دربارهٔ عناصر این سلسله‌مراتب صادق خواهد بود (بنابراین تنها \emph{یک} مجموعهٔ «تهی» می‌تواند وجود داشته باشد). اما قصد ندارم از تصور انباشتی-تکراری برای توجیه گسترش‌مندی استفاده کنم (بنگرید به~\cite{Boolos1971}). باز هم پیشنهاد من این است که گسترش‌مندی را صرفاً از آن رو بپذیریم که به گردایه‌های گسترشی علاقه‌مندیم.

\end{document}
